\input{./._relpath-to-latexroot.ltx}
\documentclass[\latexroot/\projectname]{subfiles}
\input{\latexroot/@resources/subfile-setup.ltx}

% Model.tex-specific: Redefine verbatimwrite to be a no-op in standalone mode
% (whenstandalone removed - redundant in preamble)
\makeatletter
\renewenvironment{verbatimwrite}[1]{\comment}{\endcomment}
\makeatother

\begin{document}

% When compiled as a subfile of an integrated document, path reset
% is needed so that paths are relative to the root file location
\whenintegrated{%
  \input{\latexroot/._relpath-to-latexroot.ltx}%  % Resets \latexroot and reloads all paths
}

\section{High-Frequency Identification of MP, RN, and CBI Shocks}\whenintegrated{\label{model}}
\whenintegrated{\label{sec:model}} % integrated doc: SST for labels

This section outlines the construction of target-rate and FG surprises, as well as the identification strategy used to isolate MP, RN, and CBI shocks embedded in these measures. Our approach closely follows the framework of \hyperlink{Jarociński and Karadi (2025)}{Jarociński and Karadi (2025)} and proceeds in three stages. First, we construct target-rate and FG surprises using high-frequency changes in interest rate swaps around monetary policy announcements in Mexico. Second, in the spirit of \hyperlink{Bauer and Swanson (2023a, 2023b)}{Bauer and Swanson (2023a, 2023b)}, we decompose these surprises into predictable and residual components by regressing each interest rate surprise on a set of macroeconomic and financial news observable prior to the policy announcement. Finally, we exploit the high-frequency co-movement between interest rate and stock price surprises to jointly identify the shocks of interest.



\subsection{Constructing Target-Rate and FG Surprises}

Following \hyperlink{Andrade and Ferroni (2021)}{Andrade and Ferroni (2021)}, \hyperlink{Solis (2023a)}{Solís (2023a, 2023b)}, and \hyperlink{Bolhuis et al. (2024)}{Bolhuis et al. (2024)}, we construct monetary policy surprises using high-frequency changes in interest rate swaps with maturities of 1, 3, 6, 9, and 12 months around monetary policy announcements. These surprises reflect revisions in expectations regarding both the current level of the policy rate and its future path.\footnote{The Mexican swap market references the 28-day interbank rate (TIIE-28), which is closely tied to the policy rate and serves as a benchmark for short-term bank lending.} Given the lack of intraday data, high-frequency changes in swap rates are measured as close-to-close movements from the trading day preceding the announcement to the day of the announcement. This approach is standard in the literature on EMEs, where intraday financial data are typically unavailable (see \hyperlink{Bolhuis et al. (2024)}{Bolhuis et al. (2024)}, \hyperlink{Witheridge (2024)}{Witheridge (2024)}, \hyperlink{Gomes et al. (2023)}{Gomes et al. (2023)}, and \hyperlink{Aruoba et al. (2021)}{Aruoba et al. (2021)}). Our sample includes 214 scheduled monetary policy announcements over the period January~2002 to December~2024. The starting date reflects the availability of bond yield data used in the estimations below and closely coincides with the adoption of an inflation-targeting regime by the central bank of Mexico in 2001.

Following the standard approach of \hyperlink{Gürkaynak et al. (2005)}{Gürkaynak et al. (2005)}, we extract the first two principal components (PCs) of these high-frequency changes in interest rate swaps and rotate them so that the second PC has zero loading on the 1-month swap rate. After this rotation, the first factor can be interpreted as reflecting unexpected changes in the current policy rate—closely tracked by the 1-month swap rate—and we therefore label it the \emph{target-rate surprise}. The second factor captures revisions to the expected near-term policy path, that is, movements in longer-maturity swap rates that are orthogonal to changes in the current policy rate, and we refer to it as the \emph{FG surprise}. For comparability and ease of interpretation, both surprises are rescaled so that the target-rate surprise moves one-for-one with changes in the 1-month swap rate, while the FG surprise affects the 12-month swap rate by the same magnitude as the target-rate surprise.

\subsection{Decomposing Target-Rate and FG Surprises}

In this section, we decompose the target-rate and FG surprises into two components: one that is predictable ex post using publicly available information prior to the announcement, and a residual component orthogonal to this information set. This decomposition follows the logic of \hyperlink{Bauer and Swanson (2023a, 2023b)}{Bauer and Swanson (2023a, 2023b)}, who highlight that interest rate surprises may capture not only exogenous MP shocks but also unexpected responses of the central bank to the state of the economy, particularly when market expectations about the policy reaction function are imperfect.

Formally, for each policy announcement $t$, we estimate
\begin{equation}
mps_t = \alpha + X_{t-}'\beta + u_t,
\end{equation}
where $mps_t$ denotes either the target-rate or the FG surprise, $X_{t-}$ collects macroeconomic and financial variables observable before the policy announcement (as indicated by the time subscript $t-$), and $u_t$ is the residual. We define the fitted component as $\text{FIT}_t \equiv \widehat{\alpha} + X_{t-}'\widehat{\beta}$ and the residual component as $\text{RES}_t \equiv \widehat{u}_t$, where by construction $\text{RES}_t$ is orthogonal to all variables included in $X_{t-}$. The fitted component captures the portion of the surprise that is explained by the pre-announcement information set and may be interpreted as reflecting revisions in expectations about the systematic policy response to news—that is, the RN component. The residual component isolates the unpredictable part of the surprise, which is more closely related to exogenous deviations from the systematic policy rule.

Empirical studies, including  \hyperlink{Cieslak (2018)}{Cieslak (2018)}, \hyperlink{Bauer and Swanson (2023a, 2023b)}{Bauer and Swanson (2023a, 2023b)}, and \hyperlink{Swanson (2024)}{Swanson (2024)}, among others, document a range of macroeconomic and financial variables that help predict interest rate surprises ex post. We incorporate several of these variables into $X_{t-}$, while also including a small set of additional variables that may be particularly relevant for an EME such as Mexico. Given Mexico’s status as a small open economy with strong economic linkages to the United States, $X_{t-}$ includes both domestic and external variables, primarily capturing U.S. financial and macroeconomic conditions, along with global commodity price movements.\footnote{Mexico is a small open economy with strong real and financial linkages to the U.S., often viewed as the “center country” in studies of small open economies. Roughly 80\% of Mexican exports are directed to the U.S., and a large share of foreign direct investment originates there. Prior work documents close trade and financial integration between the two economies, including substantial co-movement in term premia and short-term interest rates [see, for instance, \hyperlink{Fink (2015)}{Fink and Schüler (2015)} and \hyperlink{Carrillo et al. (2020)}{Carrillo et al. (2020)}].}


To capture domestic financial conditions, we include the log change in the nominal peso--U.S.\ dollar exchange rate and changes in the volatility implied by one-month options on the Mexican peso, both measured from three months prior to the announcement to the day before the announcement.\footnote{The volatility implied by one-month options on the Mexican peso reflects market expectations about near-term exchange rate uncertainty. This measure is widely followed by market participants and is regularly discussed in the Central Bank of Mexico’s quarterly inflation reports.} In addition, we consider three-month changes in the interbank funding rate and in the two-year Mexican government bond yield.\footnote{The interbank funding rate captures short-term conditions in the domestic money market and is closely tied to the monetary policy rate.}  To capture domestic macroeconomic conditions, we include log changes in the Overall Indicator of Economic Activity (IGAE, by its Spanish acronym) and in the core Consumer Price Index (CPI), measured from six months prior to the most recent data release available before the policy announcement.\footnotemark
\footnotetext{The IGAE follows the methodology and conceptual framework of the national accounts. The correlation between IGAE and GDP, both measured in quarterly percentage changes, is 0.99. The IGAE is subject to revisions, so the data available at a given point in time may differ from the final revised values released by the statistical office. Although the use of real-time data would be desirable, such data are unavailable for Mexican output. We therefore rely on revised data in our estimations.}\textsuperscript{,}\footnotemark
\footnotetext{ IGAE is released with a delay of approximately seven to eight weeks (approximated by 42 trading days), while consumer prices are typically released with a lag of about two weeks (approximated by 10 trading days).}

As external financial variables, we include log changes in the Standard \& Poor’s 500 stock index, changes in two-year and ten-year U.S.\ Treasury yields, and the log change in Brent crude oil prices, all measured over a three-month horizon prior to the policy announcement. For external macroeconomic conditions, we consider log changes in U.S.\ industrial production and in the U.S.\ Consumer Price Index, measured from six months prior to the most recent data release available before the policy announcement. Finally, following \hyperlink{Miranda-Agrippino and Ricco (2021)}{Miranda-Agrippino and Ricco (2021)} and \hyperlink{Swanson (2024)}{Swanson (2024)}, we include one lag of the dependent variable to account for potential serial correlation in monetary policy surprises.


\hyperlink{Table 1}{Table~1} reports the results from projecting target-rate and FG surprises onto the set of macroeconomic and financial variables described above, using the full sample from January 2002 to December 2024.\footnote{Monetary policy surprises, changes in interest rates, and variations in implied exchange rate volatility are expressed in basis points. Changes in real activity, consumer prices, oil prices, asset prices, and the exchange rate are measured as log differences multiplied by 100. Three-month windows correspond to 65 trading days. U.S.\ consumer prices and U.S.\ industrial production are typically released with a lag of about two weeks.} Results for target-rate surprises are reported in the first column, while those for FG surprises are shown in the second column. Standard errors (in parentheses) are computed using heteroskedasticity- and autocorrelation-consistent (HAC) estimators. As can be seen, the set of regressors explains a non-negligible share of the variation in monetary policy surprises, with $R^2$ values of approximately 22\% for target-rate surprises and 15\% for FG surprises, broadly in line with estimates reported in the related literature.\footnotemark
\footnotetext{The null hypothesis of no predictability is rejected at the 1\% significance level for target-rate surprises and at the 10\% level for FG surprises, based on heteroskedasticity- and autocorrelation-robust Wald tests.}\textsuperscript{,}\footnotemark
\footnotetext{As robustness exercises, we conduct two additional checks. First, we estimate an extended specification that augments the baseline regression with a broader set of controls commonly used in the literature. The resulting impulse responses based on the fitted and residual components remain largely unchanged (see \hyperlink{Table B.1}{Table~B.1} and \hyperlink{Figure B.1}{Figures~B.1 and B.2} in \hyperlink{Appendix B}{Appendix~B}). Second, we implement LASSO-based projections to allow the data to select the most relevant predictors (see \hyperlink{Table B.2}{Table~B.2} and \hyperlink{Figure B.3}{Figures~B.3 and B.4} in \hyperlink{Appendix B}{Appendix~B}). For target-rate surprises, the selected variables broadly coincide with those in the baseline information set, whereas for FG surprises the procedure selects a substantially smaller subset. The corresponding impulse responses remain broadly consistent with the baseline results, although the separation between MP and RN shocks becomes less pronounced, particularly for target-rate surprises, resulting in attenuated impulse responses to MP shocks. Overall, we interpret these results as suggesting that the baseline specification is sufficiently rich to capture an important share of the systematic component of monetary policy surprises that could otherwise introduce bias into our structural estimates if not properly separated.}


\input{Tables/Table_1_Mex}

These results indicate that interest rate surprises are correlated with publicly available economic information prior to policy announcements. While the underlying mechanisms are beyond the scope of this exercise, this pattern is consistent with the possibility that market expectations about the policy reaction function are imperfect. For the VAR analysis, however, the precise interpretation of this correlation is of secondary importance, as emphasized by \hyperlink{Swanson (2024)}{Swanson (2024)}. What matters is the empirical regularity that interest rate surprises are not fully orthogonal to pre-announcement economic and financial conditions. As we show below, failing to account for this feature leads to impulse response functions that display puzzling or potentially misleading dynamics following an MP shock.

\subsection{Disentangling MP, RN, and CBI Shocks embedded in Target-Rate and FG Surprises}

Building on \hyperlink{Jarociński and Karadi (2025)}{Jarociński and Karadi (2025)}, we exploit the high-frequency co-movement between interest rate and stock price surprises to identify MP, RN, and CBI shocks using sign and magnitude restrictions. The identification is implemented separately for target-rate and FG surprises. In each case, we use the corresponding fitted and residual components—denoted by $(\text{FIT}^{TR}, \text{RES}^{TR})$ for target-rate surprises and $(\text{FIT}^{FG}, \text{RES}^{FG})$ for FG surprises—together with a common stock-price surprise measure, $\text{SP}$. The latter is constructed as the log difference of the Mexican Stock Market Index (IPC) between the trading day preceding the policy announcement and the day of the announcement.


Our identification strategy assumes that, within a daily window around policy announcements, interest rate and equity-price surprises are predominantly driven by three structural disturbances—MP, RN, and CBI shocks. This assumption is stronger than in studies for AEs that rely on intraday windows surrounding policy announcements. Nevertheless, intraday financial data are typically unavailable for EMEs, and much of the literature therefore relies on daily windows for this type of analysis.\footnote{One exception is \hyperlink{Solis (2023a)}{Solís (2023a, 2023b)}, who employs intraday data on interest rate swaps for Mexico. However, the available dataset begins only in 2011 and lacks intraday observations for several of the interest rate swaps considered in our analysis, as well as for equity prices, which are central to our identification strategy. As emphasized by \hyperlink{Solis (2023a)}{Solís (2023a, 2023b)}, measurement error in daily swap-rate changes around policy announcements appears to be limited, as the correlation between intraday and daily changes in swap rates is close to one. To alleviate concerns that the equity-price surprise may partly reflect global financial developments unrelated to domestic monetary policy news, we orthogonalize our equity-price surprise with respect to contemporaneous changes in the S\&P~500 index, the VIX index, and the two-year Treasury yield as a robustness check. The results of this regression are reported in \hyperlink{Table B.3}{Table~B.3} in \hyperlink{Appendix B}{Appendix~B}. The corresponding impulse responses obtained using the orthogonalized measure are very similar to those reported in the baseline specification (see \hyperlink{Figure B.5}{Figures~B.5 and B.6} in \hyperlink{Appendix B}{Appendix~B}).
} In addition, we assume that each disturbance generates a distinct pattern of co-movement between interest rate and equity-price surprises.





A \emph{\textbf{pure MP shock}} is characterized by a negative co-movement between the residual component of the interest rate surprise—either $\text{RES}^{TR}$ or $\text{RES}^{FG}$—and the stock price surprise. This pattern is consistent with an unexpected policy tightening that raises discount rates and lowers expected future cash flows, leading to an immediate decline in equity prices. When identification is based on the decomposition of the target-rate surprise, we generally refer to the resulting disturbance as a \emph{pure target-rate shock}; when it is based on the decomposition of the FG surprise, we generally refer to it as a \emph{pure FG shock}.

An \emph{\textbf{RN shock}} is identified by a negative co-movement between the fitted component of the interest rate surprise—either $\text{FIT}^{TR}$ or $\text{FIT}^{FG}$—and the stock price surprise. This shock captures episodes in which the central bank responds more aggressively than markets anticipated to publicly available information that was already incorporated into expectations prior to the announcement. As a result, the main information revealed within the announcement window is the stronger-than-expected monetary policy response, which depresses equity prices through standard discount-rate and cash-flow channels.

A \emph{\textbf{CBI shock}} is identified by a positive co-movement between the overall monetary policy surprise—defined as $\text{FIT}^{j}+\text{RES}^{j}$—and the equity-price surprise. This pattern may arise when the central bank reveals favorable new information about economic fundamentals while simultaneously tightening policy. In such cases, equity prices increase in response to improved growth prospects and higher expected future cash flows, despite the policy tightening.

These assumptions are summarized by the following impact representation:
\begin{equation}
\begin{pmatrix}
\text{RES}^{j} \\
\text{FIT}^{j} \\
\text{SP}
\end{pmatrix}
=
\begin{pmatrix}
+ & c_{12} & c_{13} \\
c_{21} & + & c_{23} \\
- & - & +
\end{pmatrix}
\begin{pmatrix}
\text{MP}^{j} \\
\text{RN}^{j} \\
\text{CBI}^{j}
\end{pmatrix},
\quad j \in \{\text{TR}, \text{FG}\},
\end{equation}
together with the restriction that the overall policy surprise (the sum of fitted and residual components) rises on impact after a favorable CBI shock:
\begin{equation}
\label{eq:cbi_policy_surprise}
(c_{13}+c_{23}) > 0.
\end{equation}

Finally, we impose magnitude restrictions to ensure that the FIT--RES decomposition provides informative proxies for the MP and RN shocks:
\begin{equation}
\label{eq:mag_res_mp}
|c_{12}| < c_{11},
\end{equation}
\begin{equation}
\label{eq:mag_res_rn}
|c_{21}| < c_{22}.
\end{equation}
Restriction \eqref{eq:mag_res_mp} ensures that the residual component loads more strongly on the MP shock than on the RN shock, while restriction \eqref{eq:mag_res_rn} ensures that the fitted component loads more strongly on the RN shock than on the MP shock. 


\end{document}
